\subsection{Cập nhật repo theo nhánh \texttt{dev/NguyenDucManh} tại commit \texttt{aadd7e9}}
\label{sec:repo-update}
Tại lần đối chiếu ngày 17/11/2026, nhánh \texttt{dev/NguyenDucManh} đứng ở
\texttt{aadd7e96} (commit ngày 15/11 UTC), kế tiếp \texttt{3ca5f0c}. Bộ số liệu
định lượng canonical hiện hành nằm trong \texttt{outputs/metrics/duc\_10\_11}
(theo protocol \texttt{duc\_10\_11\_v1}); các file trong
\texttt{metrics/duc\_08\_11} và snapshot \texttt{58bdfd8} là dữ liệu lịch sử.
Commit mới nhất chủ yếu hoàn thiện demo web, không thay đổi các CSV canonical.

\subsubsection{Rà soát các nhánh từ xa}
GitHub liệt kê bảy nhánh. So sánh commit của mỗi nhánh với \texttt{dev} tại
\texttt{aadd7e96} cho trạng thái sau:
\begin{itemize}
    \item \texttt{dev/NguyenDucManh}: nhánh chuẩn đang được báo cáo, HEAD \texttt{aadd7e96}.
    \item \texttt{main}: đi sau \texttt{dev} 11 commit, không có commit riêng phía trước.
    \item \texttt{fiji/PhamBaoLoc}: phân nhánh, có 1 commit riêng và thiếu 2 commit của \texttt{dev}.
    \item \texttt{final}: phân nhánh, có 3 commit riêng và thiếu 1 commit của \texttt{dev}.
    \item \texttt{thucnghiem/BuiTrungDuc}: có 2 commit riêng và thiếu 51 commit của \texttt{dev}.
    \item \texttt{web/demo-canonical}: có 1 commit riêng và thiếu 5 commit của \texttt{dev}.
    \item \texttt{docs/report}: HEAD \texttt{bcfda814}; không có tổ tiên chung với \texttt{dev}, nên đây là lịch sử độc lập.
\end{itemize}
Vì vậy, số liệu và mô tả tính năng dưới đây theo nhánh \texttt{dev}; các nhánh
khác không được xem là nguồn thay thế protocol canonical nếu chưa hợp nhất.

\subsubsection{Metric Tầng 1 và protocol Tầng 2 hiện hành}
Protocol Tầng 1 chấm \texttt{PASS} chỉ khi sai số tương đối cực đại nhỏ hơn
$10^{-4}$ và Pearson lớn hơn $0.9999$. Kết quả trên 220 bản đồ gồm 35
\texttt{PASS} và 185 \texttt{FAIL}; các phép phân tích mẫu có thể nêu ngưỡng
rộng hơn để minh họa nhưng không làm thay đổi trạng thái canonical.
Tầng 2 gồm 10 ảnh BSDS, 10 DRIVE và 5 Blob. Bảng tổng hợp được trích từ
\texttt{KET\_QUA\_METRIC.md} và các CSV trong snapshot \texttt{duc\_10\_11}.

\begin{table*}[htbp]
\centering
\caption{F1 theo protocol canonical \texttt{duc\_10\_11\_v1}; micro-F1 được báo cáo riêng.}
\label{tab:repo-canonical-f1}
\begin{tabular}{lrrrr}
\toprule
Tác vụ & $N$ ảnh & Python & Fiji & $\Delta F_1$ (Fiji--Python)\\
\midrule
BSDS, Canny macro, $\sigma=1$ & 10 & 0.455418 & 0.455478 & $+0.000060$\\
BSDS, Canny macro, $\sigma=2$ & 10 & 0.477886 & 0.477886 & $0.000000$\\
BSDS, Canny macro, $\sigma=4$ & 10 & 0.441491 & 0.441491 & $0.000000$\\
BSDS, Canny macro, $\sigma=8$ & 10 & 0.343340 & 0.343340 & $0.000000$\\
DRIVE, Frangi macro & 10 & 0.477382 & 0.473863 & $-0.003519$\\
DRIVE, Frangi micro & 10 & 0.486776 & 0.483343 & $-0.003433$\\
Blob, LoG macro & 5 & 0.898418 & 0.897373 & $-0.001045$\\
Blob, LoG micro & 5 & 0.519135 & 0.514096 & $-0.005039$\\
\bottomrule
\end{tabular}
\end{table*}

BSDS dùng dung sai đĩa Euclidean 2 pixel. Các cặp ngưỡng Python được chọn
trên lưới đã công bố rồi giữ nguyên khi chấm Fiji; tại $\sigma=2$ là
$(p_{low},p_{high})=(67,91)$. Đây là đánh giá trên tập đã dùng chọn ngưỡng,
không phải phép đo test độc lập. DRIVE dùng Frangi tối màu, $\beta=0.5$,
$c=0.5\max(S)$ và cùng ngưỡng tuyệt đối đã chọn trên Python; vùng chấm là
giao FOV sau khi xói mòn theo support của các scale, bán kính lớn nhất hiện
tại 32 pixel. Protocol này khác phiên bản cũ dùng percentile 85 hoặc FOV
8 pixel nên không so trực tiếp các điểm giữa hai phiên bản.

Blob dùng LoG chuẩn hóa và bảy scale đã cấu hình; lưới ngưỡng
$[0.5,1,2,5]$ được chọn trên Python rồi giữ nguyên cho Fiji. Ghép một-một
dòi hỏi tâm cách không quá 2 pixel và sai số bán kính không quá 20\%; 15
vùng ignore được áp trước khi đếm FP. Trên \texttt{blob\_05} có 276 nhãn,
Python có 140 TP, 153 FP, 136 FN ($F_1=0.492091$); Fiji có 139 TP, 156 FP,
137 FN ($F_1=0.486865$). Bốn ảnh tổng hợp đạt F1 bằng 1 ở cả hai nguồn, vì
vậy macro-F1 cao cần được đọc cùng micro-F1 và kết quả riêng ảnh tế bào.

\subsubsection{Tập kiểm chứng FeatureJ/Fiji và Ground Truth}
Các tài liệu workflow Fiji ghi nhận 435 TIFF float32, gồm 423 bản đồ thuộc
phạm vi macro hiện hành và 12 file Edge legacy. Nhật ký có 369 dòng lệnh:
303 lệnh một đầu ra và 66 lệnh hai đầu ra; sau khi loại 12 dòng legacy, phạm
vi macro hiện hành gồm 357 lần gọi và 423 TIFF. Hessian có 14 ảnh, bốn bản
đồ mỗi ảnh, tương ứng 56 lần gọi và 112 TIFF. Bộ kiểm tra bỏ qua file đã có;
con số kiểm kê không khẳng định mọi TIFF được tạo mới trong cùng một lần chạy
\cite{xlaUpdate950,xlaRepository}.

Ground Truth BSDS gồm 5--7 bộ nhãn mỗi ảnh, tổng cộng 56 bộ nhãn ảnh; số này
không đồng nghĩa 56 người gán nhãn. DRIVE có mask mạch và FOV riêng. Bốn ảnh
Blob tổng hợp có bốn đối tượng mỗi ảnh. \texttt{blob\_05} có 276 đối tượng
đã duyệt và 15 vùng ignore; nhãn tế bào do nhóm rà soát, không phải nhãn
chuyên gia độc lập của nhà cung cấp \cite{xlaUpdate950,xlaRepository}.

\subsubsection{Đo thời gian ImageScience headless}
Bộ runtime mới đo 357 ca cho mỗi nguồn, mỗi ca có một warm-up và năm lần đo
giữ lại: tổng cộng 2,142 mẫu mỗi nguồn và 4,284 mẫu. Python dùng NumPy;
nguồn Fiji gọi API ImageJ 1.53k/ImageScience 3.0.0 trên JVM headless. Đọc ảnh,
chuyển kiểu, sao chép input và ghi file nằm ngoài vùng đo; đây không phải thời
gian chạy GUI hay toàn bộ pipeline chấm điểm. Median tỷ số
$T_{Python}/T_{ImageScience}$ trên 357 ca là 10.668; tỷ số lớn hơn 1 nghĩa là
Python mất nhiều thời gian hơn trong ca đó. IQR biểu đồ là phân bố giữa các
median ca, không phải khoảng tin cậy. Không khái quát kết quả sang máy hoặc
tác vụ khác. Bộ benchmark 36 ca cũ trong
\texttt{canonical\_58bdfd8/optimization\_benchmark.csv} được giữ như dữ liệu
lịch sử và không gộp với runtime mới.

\subsubsection{Demo web bốn khung}
Commit \texttt{aadd7e9} hoàn thiện demo Flask tại \texttt{src/web/}, hiển thị
ảnh gốc, kết quả Python, bản đồ tham chiếu Fiji và Ground Truth/overlay sai số.
Preset Canny dùng \texttt{bsds\_01}, $\sigma=2$, $(67,91)$; Frangi dùng
\texttt{drive\_08}, $\beta=0.5$, ngưỡng 0.05, bốn scale $[1,2,4,8]$ và FOV;
Blob dùng \texttt{blob\_05}, ngưỡng 5.0, cực tính bright và đủ bảy scale.

Tên panel được phân biệt với detector để tránh hiểu nhầm: panel Canny của Fiji
là gradient $|\nabla L|$ đã chuẩn hóa theo $\sigma$, không phải bản đồ Canny;
nếu thiếu file gradient, ứng dụng dựng từ cặp $L_x,L_y$ của Fiji và không
fallback sang file Edge. So sánh hướng dùng sai số góc quấn $2\pi$, không dùng
Pearson tuyến tính. Ảnh Python và Fiji dùng chung dải hiển thị. Panel Frangi
là chuỗi Fiji Hessian rồi Python Frangi (hybrid); panel Blob là đáp ứng cực
đại $|LoG|$ của Fiji, không phải danh sách blob. Mã mới bổ sung kiểm tra đầu
vào, bộ hủy request cũ và request ID tuần tự để tránh kết quả cũ ghi đè thao
tác mới. Demo phục vụ xem và đối chiếu; các số benchmark vẫn lấy từ CSV
canonical, không suy ra từ giao diện.

Có thể chạy theo hướng dẫn \texttt{src/web/README.md} bằng
\texttt{pip install -r requirements.txt}, sau đó
\texttt{python -m src.web.app --port 5000}. Các phụ thuộc web mới là Flask và
flask-cors.

\subsubsection{Tiêu chí nghiệm thu Tầng 1}
Kết quả canonical giữ nguyên chuẩn gốc: 35/220 bản đồ đạt đồng thời
$\mathrm{rel\_err}<10^{-4}$ và Pearson $>0.9999$; 185/220 không đạt cả hai
điều kiện. Nhãn \texttt{RELAXED} trong các bảng chạy mẫu chỉ mô tả ngưỡng
chẩn đoán rộng hơn, không thay thế \texttt{FAIL} hoặc chứng nhận nghiệm thu.
Tương quan cao không đủ để kết luận sai số biên độ đạt chuẩn.
